\section{Java Collection}

Java Collection membahas framework struktur data standar Java untuk menyimpan, mencari, mengurutkan, mengelompokkan, dan memproses object. Collection Framework memakai interface seperti \texttt{Collection}, \texttt{List}, \texttt{Set}, \texttt{Queue}, dan \texttt{Map}, serta implementasi konkret seperti \texttt{ArrayList}, \texttt{LinkedList}, \texttt{HashSet}, \texttt{TreeSet}, \texttt{HashMap}, dan \texttt{TreeMap}.

\begin{cmt}{Keterbatasan Sumber}{}
Section ini disusun berdasarkan pengetahuan umum Java. Fokus dipilih untuk kebutuhan ujian: pemilihan struktur data, kontrak equals-hashCode, ordering, generics, iterator, dan kompleksitas umum.
\end{cmt}

\subsection{Outline Fundamental Konsep Materi}

\begin{enumerate}
    \item Hierarki Collection Framework \textbf{[SUDAH DIKUASAI]}
    \begin{enumerate}
        \item \texttt{Iterable}, \texttt{Collection}, \texttt{List}, \texttt{Set}, \texttt{Queue}.
        \item \texttt{Map} bukan subtype \texttt{Collection}, tetapi bagian dari framework.
    \end{enumerate}
    \item List \textbf{[SUDAH DIKUASAI]}
    \begin{enumerate}
        \item Ordered, boleh duplikat, akses berdasarkan indeks.
        \item \texttt{ArrayList} vs \texttt{LinkedList}.
    \end{enumerate}
    \item Set \textbf{[SUDAH DIKUASAI]}
    \begin{enumerate}
        \item Tidak menyimpan duplikat.
        \item \texttt{HashSet}, \texttt{LinkedHashSet}, \texttt{TreeSet}.
    \end{enumerate}
    \item Map \textbf{[SUDAH DIKUASAI]}
    \begin{enumerate}
        \item Pasangan key-value.
        \item Key unik; value boleh duplikat.
        \item \texttt{HashMap}, \texttt{LinkedHashMap}, \texttt{TreeMap}.
    \end{enumerate}
    \item Iterator dan enhanced for
    \begin{enumerate}
        \item Iterasi aman dengan \texttt{Iterator}.
        \item \texttt{ConcurrentModificationException} saat modifikasi tidak tepat.
    \end{enumerate}
    \item Equality, hashing, dan ordering
    \begin{enumerate}
        \item \texttt{equals} dan \texttt{hashCode} untuk hash-based collection.
        \item \texttt{Comparable} dan \texttt{Comparator} untuk ordering.
    \end{enumerate}
    \item Pemilihan collection
    \begin{enumerate}
        \item Berdasarkan duplikasi, urutan, pencarian, indeks, dan sorting.
    \end{enumerate}
\end{enumerate}

\subsection{Pentingnya Materi dan Relevansinya}

\begin{jawab}[Inti Relevansi.]
Collection adalah API inti untuk menyimpan banyak object. Kemampuan memilih collection yang tepat menentukan kebenaran, performa, dan kejelasan desain program Java.
\end{jawab}

Jika data membutuhkan urutan input dan akses indeks, \texttt{List} cocok. Jika data harus unik, \texttt{Set} cocok. Jika data dicari berdasarkan key, \texttt{Map} cocok. Kesalahan memilih collection sering menyebabkan kode tambahan yang tidak perlu, performa buruk, atau bug karena duplikasi dan urutan tidak dipahami.

\subsubsection{Dependensi dan Hubungan dengan Materi Lain}

Collection bergantung pada Generics dan Interface. String sering diproses dalam collection. Stream API adalah cara deklaratif memproses collection. Equality/hashing berkaitan dengan desain class dan enkapsulasi. Wildcard muncul pada method collection seperti \texttt{addAll(Collection<? extends E>)}.

\subsection{Penjelasan Konsep Fundamental}

\subsubsection{List}

\begin{jawab}[Inti Konsep.]
\texttt{List} adalah collection berurutan yang mengizinkan duplikat dan menyediakan akses berdasarkan indeks.
\end{jawab}

\begin{lstlisting}[language=Java, caption={Penggunaan List}, label={lst:list-basic}]
List<String> names = new ArrayList<>();
names.add("Ayu");
names.add("Budi");
names.add("Ayu");

System.out.println(names.get(0)); // Ayu
System.out.println(names.size()); // 3
\end{lstlisting}

\texttt{ArrayList} umumnya menjadi pilihan default untuk list karena akses indeks cepat dan performa iterasi baik. \texttt{LinkedList} jarang menjadi pilihan default; meskipun insertion di tengah secara teori dapat murah jika node sudah diketahui, pencarian posisi tetap mahal.

\subsubsection{Set}

\begin{jawab}[Inti Konsep.]
\texttt{Set} menyimpan elemen unik. Keunikan bergantung pada \texttt{equals}/\texttt{hashCode} untuk hash set atau ordering untuk tree set.
\end{jawab}

\begin{lstlisting}[language=Java, caption={Penggunaan Set untuk menghapus duplikasi}, label={lst:set-basic}]
Set<String> uniqueNames = new HashSet<>();
uniqueNames.add("Ayu");
uniqueNames.add("Budi");
uniqueNames.add("Ayu");

System.out.println(uniqueNames.size()); // 2
\end{lstlisting}

\texttt{HashSet} tidak menjamin urutan. \texttt{LinkedHashSet} mempertahankan urutan insertion. \texttt{TreeSet} menjaga elemen terurut berdasarkan \texttt{Comparable} atau \texttt{Comparator}.

\subsubsection{Map}

\begin{jawab}[Inti Konsep.]
\texttt{Map} menyimpan pasangan key-value. Key harus unik; value boleh sama. Map cocok ketika data perlu dicari berdasarkan identitas tertentu.
\end{jawab}

\begin{lstlisting}[language=Java, caption={Penggunaan Map}, label={lst:map-basic}]
Map<String, Integer> scores = new HashMap<>();
scores.put("Ayu", 90);
scores.put("Budi", 85);
scores.put("Ayu", 95); // Mengganti value untuk key Ayu.

System.out.println(scores.get("Ayu")); // 95
\end{lstlisting}

Iterasi map biasanya dilakukan melalui \texttt{entrySet} agar key dan value tersedia bersamaan.

\begin{lstlisting}[language=Java, caption={Iterasi Map dengan entrySet}, label={lst:map-entryset}]
for (Map.Entry<String, Integer> entry : scores.entrySet()) {
    System.out.println(entry.getKey() + " = " + entry.getValue());
}
\end{lstlisting}

\subsubsection{Equals dan HashCode}

\begin{jawab}[Inti Konsep.]
Object yang digunakan sebagai key \texttt{HashMap} atau elemen \texttt{HashSet} harus memiliki kontrak \texttt{equals} dan \texttt{hashCode} yang konsisten.
\end{jawab}

Kontrak pentingnya: jika \texttt{a.equals(b)} bernilai \texttt{true}, maka \texttt{a.hashCode()} dan \texttt{b.hashCode()} harus sama. Jika kontrak ini dilanggar, lookup pada \texttt{HashMap} atau \texttt{HashSet} dapat gagal.

\begin{lstlisting}[language=Java, caption={Record otomatis memiliki equals dan hashCode berbasis komponen}, label={lst:record-map-key}]
public record StudentId(String value) { }

Map<StudentId, String> students = new HashMap<>();
students.put(new StudentId("13524044"), "Narendra");

System.out.println(students.get(new StudentId("13524044"))); // Narendra
\end{lstlisting}

\subsubsection{Comparable dan Comparator}

\begin{jawab}[Inti Konsep.]
\texttt{Comparable} mendefinisikan natural ordering pada class, sedangkan \texttt{Comparator} mendefinisikan ordering eksternal yang dapat berubah sesuai kebutuhan.
\end{jawab}

\begin{lstlisting}[language=Java, caption={Sorting dengan Comparator}, label={lst:comparator}]
List<Student> students = new ArrayList<>(List.of(
    new Student("Ayu", 90),
    new Student("Budi", 85)
));

students.sort(Comparator
    .comparingInt(Student::score)
    .reversed()
    .thenComparing(Student::name));
\end{lstlisting}

\subsection{Komponen Kunci Penting}

\begin{table}[H]
\centering
\begin{tabularx}{\textwidth}{|L{0.28\textwidth}|X|}
\hline
\textbf{Komponen Kunci} & \textbf{Makna atau Peran} \\
\hline
\texttt{List} & Collection berurutan, boleh duplikat, akses indeks. \\
\hline
\texttt{Set} & Collection unik tanpa duplikasi. \\
\hline
\texttt{Map} & Struktur key-value dengan key unik. \\
\hline
\texttt{ArrayList} & Implementasi list berbasis array dinamis. \\
\hline
\texttt{HashSet} & Set berbasis hashing. \\
\hline
\texttt{HashMap} & Map berbasis hashing. \\
\hline
\texttt{TreeSet}/\texttt{TreeMap} & Collection/map terurut. \\
\hline
\texttt{Iterator} & Object untuk mengiterasi collection. \\
\hline
\texttt{Comparable} & Natural ordering. \\
\hline
\texttt{Comparator} & Ordering eksternal. \\
\hline
\end{tabularx}
\caption{Komponen kunci dalam materi Java Collection}
\label{tab:komponen-kunci-java-collection}
\end{table}

\subsection{Algoritma dan Rumus Penting}

\subsubsection{Prosedur Memilih Collection}

\begin{jawab}[Tujuan Algoritma.]
Prosedur ini membantu memilih struktur data Java berdasarkan kebutuhan data.
\end{jawab}

\begin{lstlisting}[caption={Pseudocode pemilihan collection}, label={lst:collection-choice}]
Input: Kebutuhan penyimpanan data
Output: Jenis collection yang sesuai

1. Jika data dipetakan dari key ke value:
      gunakan Map.
2. Jika data harus unik:
      gunakan Set.
3. Jika data perlu urutan dan akses indeks:
      gunakan List.
4. Jika perlu urutan insertion pada set/map:
      gunakan LinkedHashSet atau LinkedHashMap.
5. Jika perlu urutan sorted:
      gunakan TreeSet atau TreeMap, atau sort List dengan Comparator.
6. Jika operasi utama adalah FIFO:
      gunakan Queue atau Deque.
7. Jika operasi banyak dilakukan secara concurrent:
      pertimbangkan collection dari java.util.concurrent.
\end{lstlisting}

\subsubsection{Kompleksitas Umum}

\begin{jawab}[Makna Rumus.]
Kompleksitas membantu membandingkan biaya operasi collection. Nilai berikut adalah aturan umum, bukan jaminan absolut untuk semua kondisi.
\end{jawab}

\begin{table}[H]
\centering
\begin{tabularx}{\textwidth}{|L{0.25\textwidth}|L{0.25\textwidth}|X|}
\hline
\textbf{Struktur} & \textbf{Operasi Umum} & \textbf{Kompleksitas Rata-Rata} \\
\hline
\texttt{ArrayList} & \texttt{get(i)} & $O(1)$ \\
\hline
\texttt{ArrayList} & tambah di akhir & amortized $O(1)$ \\
\hline
\texttt{HashSet}/\texttt{HashMap} & add/get/contains & rata-rata $O(1)$ \\
\hline
\texttt{TreeSet}/\texttt{TreeMap} & add/get/contains & $O(\log n)$ \\
\hline
\end{tabularx}
\caption{Kompleksitas umum collection Java}
\label{tab:kompleksitas-collection}
\end{table}

\subsection{Checklist Pemahaman}

\subsubsection{Submateri yang Telah Dibahas}

\begin{itemize}
    \item[$\square$] Saya dapat membedakan \texttt{List}, \texttt{Set}, dan \texttt{Map}.
    \item[$\square$] Saya dapat memilih \texttt{ArrayList}, \texttt{HashSet}, \texttt{HashMap}, atau struktur lain sesuai kebutuhan.
    \item[$\square$] Saya dapat menjelaskan kontrak \texttt{equals} dan \texttt{hashCode}.
    \item[$\square$] Saya dapat memakai \texttt{Comparable} dan \texttt{Comparator}.
    \item[$\square$] Saya dapat membaca kompleksitas umum collection.
\end{itemize}

\subsubsection{Prioritas Belajar}

\begin{tabularx}{\textwidth}{|L{0.22\textwidth}|X|}
\hline
\textbf{Prioritas} & \textbf{Materi yang Perlu Dikuasai} \\
\hline
\textbf{Wajib Dikuasai} & \texttt{List}, \texttt{Set}, \texttt{Map}, implementasi utama, equals-hashCode, Comparator, pemilihan collection. \\
\hline
\textbf{Cukup Paham Konsep} & Iterator, fail-fast behavior, kompleksitas umum, concurrent collection. \\
\hline
\textbf{Sudah Dikuasai} & Termasuk materi kuis, tetapi penting sebagai basis Stream API dan soal implementasi Java. \\
\hline
\end{tabularx}
